% ===================================================================
%  పట్టిక సంఖ్య 1 — పాఠశాల, ఉన్నత పాఠశాల, తరగతి గది.
%
%  Ceci n'est PAS une transcription : c'est une traduction, faite en
%  2026, en telougou d'aujourd'hui. D'ou trois differences avec
%  texto/io et texto/fr, et trois seulement :
%
%   * pas de \nl ni de \cc — la traduction ne suit aucune ligne
%     imprimee, et les lignes de ce fichier ne sont que du confort ;
%   * pas de \begin{VUpage} — elle n'a ni feuillet ni folio ;
%   * le gras se pose sur le TERME seul, ponctuation dehors.
%
%  Les cles « %%K » sont celles de texto/io, macro pour macro pour
%  l'apparat, et les renvois \textsuperscript{(n)} visent les MEMES
%  objets de la planche, dans le MEME ordre.
%
%  LA PREMIERE COLONNE DRAVIDIENNE DU RELEVE, ET C'EST UN CHANGEMENT
%  DE FAMILLE, PAS DE LANGUE. Les quatre colonnes indiennes deja
%  faites — hindi, bengali, pendjabi, marathi — sont indo-europeennes
%  et cousines de l'ido de tres loin. Le telougou ne l'est pas du
%  tout, et cela se voit a chaque phrase :
%
%    * il AGGLUTINE ses cas au lieu de les postposer : బల్ల « la
%      table », బల్లపై « sur la table », బల్లమీది « qui est sur la
%      table » — un seul mot la ou l'ido en met deux ou trois ;
%    * il n'a PAS de genre grammatical mais une opposition
%      HUMAIN / NON-HUMAIN : le pluriel de « eleve » et celui de
%      « banc » ne se construisent pas pareil, et le verbe s'accorde
%      differemment ;
%    * son VERBE FINIT LA PHRASE, sans exception et sans recours.
%
%  LA POSTPOSITION MORD DONC ENCORE PLUS QUE POUR LE MARATHI, et le
%  remede est le meme, eprouve sur seize tableaux : nommer d'abord ce
%  que l'ido nomme d'abord, puis qualifier par une seconde
%  proposition. On verra tableau par tableau ce que cela coute.
%
%  L'APPARAT CHANGE DE MOT, comme il avait change pour le marathi :
%  le telougou dit పట్టిక pour un tableau, శ్రేణి pour une serie, et
%  రంగం pour une scene — le mot de son theatre, comme le marathi
%  disait प्रवेश. Les trois motifs de html.py ont chacun gagne un
%  membre.
%
%  LA CLASSE TELOUGOUE A SON VOCABULAIRE ENTIER, ET IL N'EST NI DU
%  SANSKRIT DE PARADE NI DE L'ANGLAIS :
%
%    నల్లబల్ల (1)    le tableau noir — « la planche noire »
%    సుద్దముక్క (3)   la craie
%    పలక (42)       l'ardoise
%    బలపం           le crayon d'ardoise
%    సిరాబుడ్డి (25)  l'encrier
%    పీల్చు కాగితం (41) le buvard
%    నిఘంటువు (30)   le dictionnaire
%    వాచకం          le livre de lecture
%    చిత్తు పుస్తకం (35) le cahier de brouillon
%    కొలబద్ద (50)    la regle — « la baguette a mesurer »
%    వృత్తలేఖిని (26) le compas — « ce qui trace le cercle »
%    బడి సంచి (57)   le cartable
%    అల్మారా (15)    l'armoire
%    గవాక్షం (78)    le vasistas
%    ఆవరణ (75)      la cour
%    విరామ సమయం (76) la recreation
%    ప్రధానోపాధ్యాయుడు (71) le proviseur
%    బంట్రోతు (77)   l'homme de service
%
%  « బంట్రోతు » EST LE MEME EMPLOI QUE LE فرّاش EGYPTIEN ET LE शिपाई
%  MARATHE : l'homme de service de l'ecole, celui qui porte le
%  registre des absents. Troisieme colonne d'affilee, trois ecoles
%  reelles, et la meme fonction nommee sans periphrase.
%
%  « postigo » (renvoi 78) EST UN VASISTAS, et le telougou a le mot :
%  గవాక్షం, la lucarne haute qu'on ouvre pour l'air. Meme resultat
%  que dans les trente-quatre autres colonnes.
%
%  LA GAMME RESTE OCCIDENTALE, ET C'EST VOULU. Le telougou a sa
%  propre solmisation — సరిగమపదని — vieille de plus de mille ans, et
%  l'ecrire ici serait tentant. Mais le tableau nomme do-re-mi et
%  pose l'equivalence « = C. D. E. F. G. A. B. » : c'est la gamme
%  d'Europe qui est l'objet, et l'on ne transpose pas les choses. En
%  revanche « తాళం వేయు », battre la mesure (68), est du telougou
%  ordinaire, parce que le tala est ce qu'on bat ici aussi.
% ===================================================================

%%K t01-apar-0 apar
\VUsaut{2.0mm}
\VUcentre{12.6pt}{120}{వివరణ - పుస్తిక}
\VUsaut{3.2mm}
\VUcentre{8.4pt}{160}{యొక్క}
\VUsaut{2.6mm}
\VUtitre{15.6pt}{0}{డెల్మా - పట్టికలకు తోడ్పడే}
\VUsaut{3.4mm}
\VUornamento{9pt}
\VUsaut{5.0mm}
\VUcentre{13.2pt}{90}{మొదటి శ్రేణి}
\VUsaut{3.0mm}
\VUfilet{16mm}
\VUsaut{5.6mm}

%%K t01-apar-1 apar
\VUtitre{15.0pt}{60}{పట్టిక సంఖ్య 1}
\VUsaut{1.4mm}
\VUtitre{11.4pt}{0}{పాఠశాల, ఉన్నత పాఠశాల, తరగతి గది.}
\VUsaut{2.2mm}
\VUfilet{20mm}
\VUsaut{6.4mm}

%%K t01-tit-1 sub
\VUpk{10.2pt}{40}{సాధారణ వర్ణన.}
\VUsaut{3.0mm}

%%K t01-01-1 p
1. --- ఈ పట్టిక ఒక \VUgras{ఉన్నత పాఠశాలలోని} \VUgras{తరగతి గదిని} చూపిస్తుంది. ఈ తరగతి గదిలో
ఎనిమిది \VUgras{బల్లలు} \textsuperscript{(18)} మరియు ఎనిమిది \VUgras{బెంచీలు}
\textsuperscript{(32)} ఉన్నాయి. ప్రతి బెంచీపై ముగ్గురు విద్యార్థులు కూర్చుంటారు.
\VUgras{ఉపాధ్యాయుడు} \textsuperscript{(7)} ఒక \VUgras{కుర్చీపై} \textsuperscript{(8)}
కూర్చుంటాడు, తన \VUgras{పీఠంలో} \textsuperscript{(6)}.

%%K t01-02-1 p
2. --- పీఠానికి ఎడమవైపు \VUgras{గాజు తలుపుల} \VUgras{అల్మారా} \textsuperscript{(15)} ఉంది.
కుడివైపు \VUgras{నల్లబల్ల} \textsuperscript{(1)}, దానితోపాటు \VUgras{తుడిచే గుడ్డ}
\textsuperscript{(2)} మరియు \VUgras{సుద్దముక్క} \textsuperscript{(3)} ఉన్నాయి.

%%K t01-02-2 p
పీఠం పైన \VUgras{గోడ పటాలు} \textsuperscript{(9, 11,} \textsuperscript{12)} మరియు ఒక
\VUgras{పటం} \textsuperscript{(10)} వేలాడుతున్నాయి.

%%K t01-03-1 p
3. --- విద్యార్థులందరూ తమ \VUgras{స్థానంలో} \textsuperscript{(17)} కూర్చోలేదు. జాన్
\textsuperscript{(4)} నల్లబల్లదగ్గర నిలబడి, \VUgras{తుడిచే గుడ్డను} పట్టుకుని
\VUgras{జ్యామితి బొమ్మలను} చెరిపేస్తున్నాడు.

%%K t01-03-2 p
పీటర్ పీఠం దగ్గర ఉన్నాడు; అతను \VUgras{రాత పనులను} \textsuperscript{(5)} సేకరించి
ఉపాధ్యాయుడికి అందిస్తున్నాడు.

%%K t01-04-1 p
4. --- \VUgras{ఈ} ఉపాధ్యాయుడు \VUgras{సజీవ భాషల} ఉపాధ్యాయుడు. విదేశీ భాషలు
\VUgras{(జర్మన్, ఇంగ్లిష్, స్పానిష్, ఇటాలియన్, రష్యన్} లేదా \VUgras{ఫ్రెంచ్)} మాట్లాడటం,
చదవడం, రాయడం అతను విద్యార్థులకు నేర్పుతాడు.

%%K t01-05-1 p
5. --- తరగతి గది చాలా విశాలంగా, చాలా వెలుతురుగా ఉంది. చివర గాలి, వెలుతురు వచ్చేందుకు రెండు
\VUgras{గవాక్షాలతో} \textsuperscript{(78)} ఒక వెడల్పైన \VUgras{కిటికీ} మరియు ఒక
\VUgras{గాజు తలుపు} ఉన్నాయి.

%%K t01-05-2 p
ఈ తరగతి గది చల్లగా ఉండదు. దాన్ని వెచ్చగా ఉంచేందుకు మధ్యలో ఒక మంచి \VUgras{పొయ్యి}
\textsuperscript{(46)} ఉంది, తన \VUgras{గొట్టంతో} \textsuperscript{(45)}.

%%K t01-05-3 p
గోడకు \VUgras{బట్టల కొక్కేలు} \textsuperscript{(58)} బిగించి ఉన్నాయి; వాటికి విద్యార్థుల
టోపీలు, కోట్లు వేలాడుతున్నాయి.

%%K t01-tit-2 sub
\VUsaut{5.2mm}
\VUpk{10.2pt}{40}{ఆట స్థలం}
\VUsaut{3.6mm}

%%K t01-06-1 p
6. --- \VUgras{గడియారం} \textsuperscript{(63)} తొమ్మిది గంటల ఐదు నిమిషాలు చూపిస్తోంది.

%%K t01-06-2 p
\VUgras{విరామ సమయం} \textsuperscript{(76)} ఇప్పుడే ముగిసింది, పాఠం మళ్ళీ మొదలవుతోంది. ఆట స్థలం \VUgras{ఆవరణలో} \textsuperscript{(75)} ఉంది. ఆడుతున్న కొందరు విద్యార్థులు మనకు కనిపిస్తున్నారు. ఒక \VUgras{సహాయ ఉపాధ్యాయుడు} \textsuperscript{(74)} వాళ్ళను కనిపెడుతున్నాడు.

%%K t01-07-1 p
7. --- ఆవరణలో ఇంకా రకరకాల మనుషులు మనకు కనిపిస్తారు — అంటే \VUgras{తనిఖీ అధికారి}
\textsuperscript{(73)}, పాఠశాల అధిపతి \VUgras{(ప్రధానోపాధ్యాయుడు)} \textsuperscript{(71)}
మరియు \VUgras{పర్యవేక్షకుడు} \textsuperscript{(72)}.

%%K t01-07-2 p
తనిఖీ అధికారి పాఠశాలలను తనిఖీ చేస్తాడు, ప్రధానోపాధ్యాయుడు పాఠశాలను నడుపుతాడు, పర్యవేక్షకుడు
చదువులను కనిపెడతాడు.

%%K t01-07-3 p
నాలుగవ వ్యక్తి \VUgras{బంట్రోతు} \textsuperscript{(77)}. గైర్హాజరులను నమోదు చేసేందుకు అతను
చంకలో ఒక రిజిస్టరు పట్టుకుని ఉన్నాడు.

%%K t01-08-1 p
8. --- ఆవరణ చివర \VUgras{వ్యాయామ సాధనాలు} \textsuperscript{(70)} మరియు \VUgras{చేతిపనుల}
\textsuperscript{(69)} కార్యశాల ఉన్నాయి.

%%K t01-08-2 p
పట్టికకు కుడివైపు \VUgras{సంగీతం} మరియు \VUgras{చిత్రలేఖనం} \VUgras{తరగతి గదులు} మనకు
కనిపించడం మొదలవుతుంది.

%%K t01-noto-1 noto
\VUnotes{143.2mm}{%
(*) \textit{నామకరణ పేర్లనూ} లాటిన్ రూపంలోనే రాయమని సలహా. లెక్కలేనన్ని భేదాలను
తప్పించుకునేందుకు అదొక్కటే మార్గం. అయినా
\textit{Giovanni, Jean, Johann, Jan, Hans, John, Ivan,} మొదలైన వాటిలో ఎంపిక ఎలా చేయాలి? పేర్ల
కోసం ఒక ప్రత్యేక నిఘంటువు తయారు చేయాలా? లాటిన్ నుంచి వాటి ప్రథమా విభక్తి రూపాన్ని తీసుకుని ఇలా
అందాం: \VUgras{Ioannes, Ioanna; Iulius, Iulia; Iakobus;} \VUgras{Andreas; Lukas.} (సంపూర్ణ
వ్యాకరణం).%
}

%%K t01-tit-3 sub
\VUpk{10.2pt}{40}{వివరమైన వర్ణన.}
\VUsaut{2.4mm}

%%K t01-tit-4 sub
\VUpk{10.2pt}{40}{మంచి విద్యార్థులు.}
\VUsaut{3.0mm}

%%K t01-09-1 p
9. --- పీఠానికి కుడివైపు మొదటి బెంచీలోని విద్యార్థులు \VUgras{హెన్రీ} (*), \VUgras{స్టీఫెన్}
మరియు \VUgras{చార్లెస్}.

%%K t01-09-2 p
హెన్రీ చాలా శ్రద్ధగా, కష్టపడి పనిచేసేవాడు; అతను ఉపాధ్యాయుడి మాట వింటాడు. తన బల్లపై అతనికి
\VUgras{నోటుపుస్తకం} \textsuperscript{(28)}, \VUgras{పుస్తకం} \textsuperscript{(29)},
\VUgras{నిఘంటువు} \textsuperscript{(30)} మరియు \VUgras{కలం పెట్టె} \textsuperscript{(31)}
ఉన్నాయి. అతని పుస్తకంలో చాలా \VUgras{పుటలు} ఉన్నాయి, ప్రతి అధ్యాయంలో కొన్ని
\VUgras{పరిచ్ఛేదాలు}, ప్రతి \VUgras{పంక్తిలో} చాలా \VUgras{పదాలు} ఉన్నాయి.

%%K t01-09-4 p
అతని పక్కనున్న స్టీఫెన్ \textsuperscript{(24)} ఉత్సాహవంతుడు. వచ్చే సారి పాఠాన్ని అతను తన
నోటుపుస్తకంలో రాసుకుంటున్నాడు. అతని \VUgras{చేతిరాత} అందంగా ఉంటుంది. అతని పుస్తకం మూసి ఉంది;
మనకు \VUgras{అట్ట} \textsuperscript{(27)} మాత్రమే కనిపిస్తుంది. అతని దగ్గరే
\VUgras{వృత్తలేఖిని} \textsuperscript{(26)} ఉంది.

%%K t01-09-5 p
చార్లెస్ \textsuperscript{(23)} తన పుస్తకాన్ని చేత్తో పట్టుకుని, పాఠాన్ని మళ్ళీ చదువుకునేందుకు
దాన్ని తెరుస్తున్నాడు. ప్రతి విద్యార్థిలాగే అతని ముందూ ఒక \VUgras{సిరాబుడ్డి}
\textsuperscript{(25)} ఉంది. అతను మాట వినేవాడు.

%%K t01-10-1 p
10. --- పక్క బల్ల దగ్గర \VUgras{లూయిస్}, \VUgras{ఆంటోని} మరియు \VUgras{పాల్} ఉన్నారు. లూయిస్
తన \VUgras{పాఠాన్ని} \textsuperscript{(22)} అప్పజెబుతూ ఉపాధ్యాయుడి \VUgras{ప్రశ్నలకు}
జవాబిస్తున్నాడు. ఆంటోని \textsuperscript{(21)} చాలా అజాగ్రత్తపరుడు; అప్పుడప్పుడూ ఉపాధ్యాయుడు
అతన్ని శిక్షిస్తాడు కూడా. పాల్ \textsuperscript{(20)} మంచి \VUgras{స్నేహితుడు} — మృదువుగా,
సాయపడేవాడిగా ఉంటాడు. అతను చాలా శ్రద్ధగలవాడు.

%%K t01-tit-5 sub
\VUsaut{4.2mm}
\VUpk{10.2pt}{40}{చెడ్డ విద్యార్థులు.}
\VUsaut{3.2mm}

%%K t01-11-1 p
11. --- హెన్రీ వెనుక \VUgras{జార్జ్} \textsuperscript{(38)} ఉన్నాడు. అతను చెడ్డవాడేమీ కాదు,
కానీ \VUgras{కబుర్లు} చెప్పేవాడు, అజాగ్రత్తపరుడు, అల్లరివాడు. అతని \VUgras{చపలత్వం} మరియు
\VUgras{అవిధేయత} పాఠం సమయంలో \VUgras{గందరగోళం} కలిగిస్తాయి, తరచూ కఠినమైన \VUgras{హెచ్చరికకు}
అతను అర్హుడవుతాడు. అతని \VUgras{చిత్తు పుస్తకం} \textsuperscript{(35)} మురికిగా ఉంది; తన
\VUgras{కలంతో} \textsuperscript{(34)} దానిపై అతను \VUgras{సిరా మరకలు} చేస్తాడు, అందుకే ఎప్పుడూ
తన \VUgras{రబ్బరును} \textsuperscript{(36)} వాడుతూ ఉంటాడు; అతని \VUgras{పుస్తకాల సంచి}
\textsuperscript{(33)} చినిగిపోయింది, ఎందుకంటే అతను చాలా అశ్రద్ధపరుడు. \VUgras{సజీవ భాషల}
తరగతి గదికి తన \VUgras{పెన్సిల్ పెట్టెను} \textsuperscript{(37)} అతను ఎందుకు తెచ్చుకుంటాడు?

%%K t01-11-2 p
అతని పక్కనున్న ఎడ్మండ్ \textsuperscript{(39)}కు అతనంటే ఇష్టం లేదు. నిజమే, ఎడ్మండ్ కొంచెం
బద్ధకస్తుడు, కానీ మౌనంగా, ప్రశాంతంగా ఉంటాడు. అతను కబుర్లు చెప్పడు; కానీ వినడానికి బదులు తన
\VUgras{వాచకంలోని} \VUgras{కథలు} చదవడానికే ఇష్టపడతాడు.

%%K t01-11-3 p
ఈ బెంచీలోని మూడో విద్యార్థి \VUgras{లూయిస్} \textsuperscript{(40)}. అతను శ్రద్ధగలవాడు. తెల్లని
\VUgras{కాగితంపై} అతను రాస్తున్నాడు. తన ముందు \VUgras{పీల్చు కాగితం} \textsuperscript{(41)}
పెట్టుకున్నాడు.

%%K t01-12-1 p
12. --- ఉపాధ్యాయుడికి ఎడమవైపు రెండో బెంచీలోని విద్యార్థులు చాలా చిన్నవాళ్ళు. మొదటివాడు, చిన్న
\VUgras{ఎమిల్}, ఇంకా బాగా రాయడం రాదు; అతను తన \VUgras{పలక} \textsuperscript{(42)}, తన
\VUgras{బలపం} మరియు తన \VUgras{స్పాంజి} \textsuperscript{(43)} తెచ్చుకుంటాడు.

%%K t01-12-2 p
అతని పక్కన \VUgras{ఆగస్ట్} మరియు \VUgras{బెర్నార్డ్} \textsuperscript{(44)} ఉన్నారు. ఈ
రెండోవాడు బాగా చదువుతాడు, కానీ అతని \VUgras{దుస్తులు} చాలా అశ్రద్ధగా ఉంటాయి.

%%K t01-13-1 p
13. --- వాళ్ళ వెనుక ముగ్గురు \VUgras{సోమరులు} ఉన్నారు. \VUgras{ఆల్బర్ట్} తన పాఠాలు నేర్చుకోడు.
\VUgras{జేకబ్} \textsuperscript{(47)} వినడానికి బదులు నిద్రపోతున్నాడు. అతని \VUgras{సోమరితనం}
అతనికి \VUgras{మందలింపులు}, \VUgras{శిక్షలు} కూడా తెచ్చిపెడుతుంది. చివరగా \VUgras{క్రిస్టియన్}
\textsuperscript{(48)} మరీ తేలికగా తీసుకుంటాడు. అతను చెడుగా \VUgras{ప్రవర్తిస్తాడు}, తనను తాను
దిద్దుకునేందుకు ఏ ప్రయత్నమూ చేయడు.

%%K t01-14-1 p
14. --- దానికి భిన్నంగా అతని పక్కవాళ్ళు మంచి నడవడిక కలవాళ్ళు. \VUgras{ఎర్నెస్ట్}
\textsuperscript{(51)}కు క్రమం, చక్కదనం ఇష్టం; తన \VUgras{కొలబద్దను} \textsuperscript{(50)}
మరియు తన \VUgras{పెన్సిల్‌ను} \textsuperscript{(49)} అతను ఎప్పుడూ వాడతాడు. అతను
\VUgras{బయటి విద్యార్థి}.

%%K t01-14-2 p
\VUgras{ఫ్రాన్సిస్} \textsuperscript{(52)} \VUgras{వసతి విద్యార్థి}; అతను ఏకరూప దుస్తులు వేసుకుని, పాఠశాలలోనే తిని, పడుకుంటాడు. అతను మంచి \VUgras{స్నేహితుడు}.

%%K t01-14-3 p
మారిస్ తన \VUgras{చాకుతో} \textsuperscript{(56)} తన \VUgras{పెన్సిల్‌ను} చెక్కుతున్నాడు; అతను
చాలా జాగ్రత్తపరుడు.

%%K t01-14-4 p
తన పుస్తకాలన్నీ అతను తెచ్చుకుంటాడు — తన \VUgras{వ్యాకరణం} \textsuperscript{(55)}, తన
\VUgras{చిన్న నోటుపుస్తకం} \textsuperscript{(54)}, చివరకు \VUgras{రాత దిండు}
\textsuperscript{(53)} కూడా. అతని \VUgras{బడి సంచి} \textsuperscript{(57)} అతని వెనుక నేలపై
ఉంది.

%%K t01-14-5 p
తరగతి గది చివరనున్న రెండు బల్లలను \VUgras{మంచి విద్యార్థులు} \textsuperscript{(19)}
ఆక్రమించుకున్నారు.

%%K t01-tit-6 sub
\VUsaut{4.6mm}
\VUpk{10.2pt}{40}{చిత్రలేఖనం మరియు సంగీతం.}
\VUsaut{3.4mm}

%%K t01-15-1 p
15. --- \VUgras{చిత్రలేఖన గదిలో} గోడకు తగిలించిన అందమైన \VUgras{నమూనాలు} ఉన్నాయి, పైకప్పుకు
వేలాడదీసిన \VUgras{దీపఛత్రంతో} ఒక \VUgras{దీపం} \textsuperscript{(62)} ఉంది.
\VUgras{ఉపాధ్యాయుడు} \textsuperscript{(60)} చాలా ఓపికగలవాడు; విద్యార్థుల \VUgras{బొమ్మలను}
\textsuperscript{(61)} అతను దిద్దుతాడు, ప్రకృతిని చూసి \VUgras{అర్ధప్రతిమను}
\textsuperscript{(59)} గీయడం వాళ్ళకు నేర్పుతాడు.

%%K t01-16-1 p
16. --- \VUgras{సంగీత గది} చాలా ఆసక్తికరంగా అనిపిస్తుంది. అక్కడ \VUgras{సంగీతం} చదవడం,
\VUgras{స్వరపంక్తులపై} రాసిన \VUgras{స్వరాలను} \textsuperscript{(67)} (do, re, mi, fa, sol,
la, si\,=\,C. D. E. F. G. A. B.) పాడటం, \VUgras{క్లెఫ్‌లను} తెలుసుకోవడం, మధురమైన
\VUgras{బాణీలు} పాడటం నేర్చుకుంటారు. సంగీత పుస్తకాలను \VUgras{పీఠంపై} \textsuperscript{(65)}
పెడతారు. విద్యార్థుల్లో ఒకడు \VUgras{పియానో వాయిస్తున్నాడు} \textsuperscript{(64)},
\VUgras{సంగీత ఉపాధ్యాయుడు} \textsuperscript{(66)} \VUgras{తాళం వేస్తున్నాడు}
\textsuperscript{(68)}.

%%K t01-17-1 p
17. --- \VUgras{చేతిపనుల} \textsuperscript{(69)} \VUgras{కార్యశాల} మూల మనకు కనిపించడం
మొదలవుతుంది; అక్కడ పిల్లలు కలపను చెక్కడం, ఇనుమును అరగదీయడం నేర్చుకోవచ్చు. ఇది అన్ని ఉన్నత
పాఠశాలల్లోనూ ఉండదు.

%%K t01-tit-7 sub
\VUsaut{5.0mm}
\VUpk{10.2pt}{40}{శాస్త్రాల తరగతి.}
\VUsaut{3.6mm}

%%K t01-18-1 p
18. --- గణిత ఉపాధ్యాయుడు విద్యార్థులకు \VUgras{జ్యామితి, అంకగణితం, లెక్కలు} మరియు
\VUgras{మెట్రిక్ పద్ధతి} నేర్పుతాడు. నల్లబల్లపై ముఖ్యమైన జ్యామితి బొమ్మలు మనకు కనిపిస్తాయి :
\VUgras{వృత్తం} \textsuperscript{(a)}, \VUgras{చతురస్రం} \textsuperscript{(b)},
\VUgras{త్రిభుజం} \textsuperscript{(c)}, \VUgras{రేఖ} \textsuperscript{(d)}.

%%K t01-18-2 p
ఒక \VUgras{పరిక్రియ} కూడా మనకు కనిపిస్తుంది; అది \VUgras{కూడిక} కాదు, \VUgras{తీసివేత} కాదు,
\VUgras{గుణకారం} కాదు — అది \VUgras{భాగహారం} \textsuperscript{(e)}.

%%K t01-19-1 p
19. --- మెట్రిక్ పద్ధతి పట్టికపై మనకు కనిపించేవి : \VUgras{మీటరు} \textsuperscript{(a)},
\VUgras{త్రాసు} \textsuperscript{(b)} మరియు దాని \VUgras{తూనికరాళ్ళు}, \VUgras{లీటరు}
\textsuperscript{(e)}, \VUgras{దశలీటరు} \textsuperscript{(f)}, \VUgras{దెసిలీటరు} మరియు
\VUgras{ఘనమీటరు} \textsuperscript{(d)}.

%%K t01-19-2 p
విద్యార్థులు \VUgras{ప్రకృతి శాస్త్రం} \textsuperscript{(11)}, జంతుశాస్త్రం
\textsuperscript{(a)}, వృక్షశాస్త్రం \textsuperscript{(b)}, భూగర్భశాస్త్రం
\textsuperscript{(c)} మరియు \VUgras{చరిత్ర} \textsuperscript{(12)} కూడా చదువుతారు.

%%K t01-19-3 p
అల్మారాలో \VUgras{భౌతిక శాస్త్రం} \textsuperscript{(15)} మరియు \VUgras{రసాయన శాస్త్ర}
\textsuperscript{(16)} సాధనాలు ఉన్నాయి.

%%K t01-19-4 p
అల్మారా పైన ఉన్నవి : \VUgras{గోళం} \textsuperscript{(14)}, \VUgras{శంకువు} మరియు
\VUgras{భూగోళం} \textsuperscript{(13)}.

%%K t01-19-5 p
పట్టికల పైన ప్రపంచంలోని ఐదు భాగాలను చూపే \VUgras{పటం} వేలాడుతోంది : \VUgras{యూరప్}
\textsuperscript{(g)}, \VUgras{ఆసియా} \textsuperscript{(e)}, \VUgras{ఆఫ్రికా}
\textsuperscript{(f)}, \VUgras{అమెరికా} \textsuperscript{(ab)} మరియు \VUgras{ఓషియానియా}
\textsuperscript{(h)}, ఇంకా రెండు పెద్ద మహాసముద్రాలు — \VUgras{పసిఫిక్} \textsuperscript{(c)}
మరియు \VUgras{అట్లాంటిక్} \textsuperscript{(d)}.

\VUsaut{6.0mm}
\VUfilet{22mm}
